\documentclass[a4paper,10pt]{book}\usepackage[]{graphicx}\usepackage[]{color}
%% maxwidth is the original width if it is less than linewidth
%% otherwise use linewidth (to make sure the graphics do not exceed the margin)
\makeatletter
\def\maxwidth{ %
  \ifdim\Gin@nat@width>\linewidth
    \linewidth
  \else
    \Gin@nat@width
  \fi
}
\makeatother

\definecolor{fgcolor}{rgb}{0.251, 0.251, 0.251}
\newcommand{\hlnum}[1]{\textcolor[rgb]{0,0.533,0.298}{#1}}%
\newcommand{\hlstr}[1]{\textcolor[rgb]{0,0.533,0.298}{#1}}%
\newcommand{\hlcom}[1]{\textcolor[rgb]{1,0.314,0.314}{#1}}%
\newcommand{\hlopt}[1]{\textcolor[rgb]{0.251,0.251,0.251}{#1}}%
\newcommand{\hlstd}[1]{\textcolor[rgb]{0.251,0.251,0.251}{#1}}%
\newcommand{\hlkwa}[1]{\textcolor[rgb]{0.502,0.627,0.188}{#1}}%
\newcommand{\hlkwb}[1]{\textcolor[rgb]{0.69,0.424,0.345}{#1}}%
\newcommand{\hlkwc}[1]{\textcolor[rgb]{0.188,0.631,0.533}{#1}}%
\newcommand{\hlkwd}[1]{\textcolor[rgb]{0.6,0,0}{#1}}%

\usepackage{framed}
\makeatletter
\newenvironment{kframe}{%
 \def\at@end@of@kframe{}%
 \ifinner\ifhmode%
  \def\at@end@of@kframe{\end{minipage}}%
  \begin{minipage}{\columnwidth}%
 \fi\fi%
 \def\FrameCommand##1{\hskip\@totalleftmargin \hskip-\fboxsep
 \colorbox{shadecolor}{##1}\hskip-\fboxsep
     % There is no \\@totalrightmargin, so:
     \hskip-\linewidth \hskip-\@totalleftmargin \hskip\columnwidth}%
 \MakeFramed {\advance\hsize-\width
   \@totalleftmargin\z@ \linewidth\hsize
   \@setminipage}}%
 {\par\unskip\endMakeFramed%
 \at@end@of@kframe}
\makeatother

\definecolor{shadecolor}{rgb}{.97, .97, .97}
\definecolor{messagecolor}{rgb}{0, 0, 0}
\definecolor{warningcolor}{rgb}{1, 0, 1}
\definecolor{errorcolor}{rgb}{1, 0, 0}
\newenvironment{knitrout}{}{} % an empty environment to be redefined in TeX

\usepackage{alltt}

\definecolor{themecolor}{RGB}{0,120,2}% % theme color

\usepackage{lipsum}
% **************************************
% Graphics
% **************************************
\usepackage{adjustbox}
\usepackage{graphicx, rotating}
%\usepackage[svgnames, table]{xcolor}% put before pstricks, textpos, tikz
%\DeclareGraphicsRule{.mps}{eps}{.mps}{}
\graphicspath{{figure/}{img/}}
% **************************************
% -> color
% **************************************
\usepackage{../../set/mycolor}
% **************************************
% <- color
% **************************************

% **************************************
% Table
% **************************************
\usepackage{booktabs,longtable,multirow,tabularx,warpcol}
\usepackage{threeparttable}
% **************************************
% Structure
% **************************************
\usepackage{../../set/mystructure}
% **************************************
% <- Structure
% **************************************

\usepackage{multind}
\makeindex{people}
\makeindex{org}
%\usepackage{titlesec}
\RequirePackage[center,pagestyles]{titlesec}

\newcommand{\company}{厦门怡惟服饰股份有限公司}             %%%%%%%%%%%%%%%%%%%%%%
\newcommand{\project}{}             %%%%%%%%%%%%%%%%%%%%%%
\newcommand{\lnotesversion}{1.0}             %%%%%%%%%%%%%%%%%%%%%
\newcommand{\lnotesdate}{2014年5月}          %%%%%%%%%%%%%%%%%%%%%%


% **************************************
% -> Languagues and Fonts
% **************************************
\usepackage{../../set/zhfonts-book}
% **************************************
% <- Languagues and Fonts
% **************************************
\usepackage{comment}
\begin{comment}
\makeatletter
\renewcommand{\tableofcontents}{%
\setlength{\columnsep}{2.5em}
%\setlength{\columnseprule}{.8pt}
\begin{multicols}{2}[\chapter*{\contentsname}]%
    \@starttoc{toc}%
\end{multicols}}
\makeatother
\end{comment}
% **************************************
% Format
% **************************************
\usepackage{calc}
\usepackage{indentfirst,setspace}

%\setlist[1]{labelindent=0.5\parindent,leftmargin=\parindent}

\usepackage{fancyvrb,listings,verbatim}

\VerbatimFootnotes
\usepackage{../../packages/ldemo}

%\usepackage{framed}
\usepackage{../../set/mymdframed}
\usepackage{marginnote}

%%
\makeatletter
 \newcommand{\cneqref}[1]{%
 方程（\ref{#1}）%
 }
\makeatother
%%
%%
\makeatletter
 \newcommand{\tabref}[1]{%
 表（\ref{#1}）%
 }
\makeatother
%%
%%
\makeatletter
 \newcommand{\figref}[1]{%
 图（\ref{#1}）%
 }
\makeatother
%%

\setlength{\abovecaptionskip}{4pt}
\setlength{\belowcaptionskip}{8pt}

%\renewcommand{\labelitemi}{\ensuremath{\bullet}}% xunicode changes the bullet

\makeatletter
\renewenvironment{quotation}{
    \list{}{
        \listparindent 2em
        \itemindent    \listparindent
        \rightmargin   \leftmargin
        \parsep        \z@ \@plus\p@
    }
    \item
}{
    \endlist
}
\makeatother

% **************************************
% -> Layout
% **************************************
\usepackage{../../set/myfancy}
\usepackage{sidenotes}

% ----- Title Style -----
\newcommand*{\base}[1]{
 \includegraphics[scale=0.3]{../../imgs/greenbase.pdf}        %%%%%%%%%%%%%%%%%%%%%%%%%%%%%%%%%%%%%%%
 }

% **************************************
% <- Layout
% **************************************

% **************************************
% Bibliography
% **************************************
\usepackage{chapterbib}
\usepackage[sectionbib,super,square,sort&compress]{natbib}

% **************************************
% -> Math
% **************************************
\usepackage{../../set/math-book}
% **************************************
% <- Math
% **************************************

\newcommand{\R}{\textsf{R}}
\makeatletter
 \newcommand{\rcode}[1]{\textsf{#1}}
\makeatother
\IfFileExists{upquote.sty}{\usepackage{upquote}}{}

\begin{document}

%<knitr: GLOBAL SETTING>========================================================




%>knitr: GLOBAL SETTING<========================================================


\frontmatter
\begin{titlepage}

\setlength\parindent{0pt}

\definecolor{mygreen}{HTML}{9ABA5A}
\definecolor{myblue}{HTML}{4F80BD}
\textblockcolour{myblue}
\textblockrulecolour{white}

% Green block%
\begin{textblock}{6}(10,0)
    \rule{0mm}{420mm}
\end{textblock}

% Three green bars%
\begin{textblock}{.02}(9.7,0)
    \rule{0mm}{420mm}
\end{textblock}

\begin{textblock}{.02}(9.8,0)
    \rule{0mm}{420mm}
\end{textblock}

\begin{textblock}{.02}(9.9,0)
    \rule{0mm}{420mm}
\end{textblock}

\begin{textblock}{3}(11,1)
    {\Huge \textbf{华昱资本}\\[5pt] \textbf{\hspace{2cm}同创价值}}
\end{textblock}

\begin{textblock}{12}(1.8,6)
\textblockcolour{}
\begin{center}
    \includegraphics[scale=1.2]{imgs/hytc.png}
\end{center}
\end{textblock}

\begin{textblock}{4}(11,12.5)
    \vspace{2cm}
    {\huge \textit{方莲}}\\[5pt]
    {\Large \lnotesdate}
\end{textblock}

\TPshowboxestrue
\setlength\TPboxrulesize{0.8pt}
\textblockcolour{themecolor}

\begin{textblock}{14}(0.01,3)
    \centering
    ~\\[20pt]
    {\fontsize{32}{40}\selectfont \CJKfamily{zhsong}{\color{white} \company\ \project\ }}\\[8pt]
    {\fontsize{36}{42} \CJKfamily{new}{\color{white}投资价值分析报告}}\\[8pt]
    版\ 本：\ \ \lnotesversion\\[20pt]
\end{textblock}
~
\end{titlepage}

\newpage
\thispagestyle{empty}
~

\begin{titlepage}
%\thispagestyle{empty}
\begin{center}
{\hrule height 2pt}\ \\[5pt]
{\fontsize{32}{40}\selectfont \CJKfamily{zhsong}{\color{themecolor} XX 公司 XX 项目}}\\[8pt]
{\fontsize{36}{42} \CJKfamily{new}{\color{themecolor}投资价值分析报告}}\\[8pt]
\normalsize 版\ 本：\  \ \lnotesversion\\[20pt]
{\hrule height 2pt}\  \\[100pt]

\huge \textit{方 \hspace{2pt} 莲}\\[8pt]
\Large \lnotesdate\\[200pt]

\textsc{HY\ Capital\ Investment\ GP}\\[8pt]
\huge \textit{华\ 昱\ 资\ 本\ 管\ 理\ 公\ 司}
\normalsize
\end{center}
\end{titlepage}
                     % Front cover and title page
\thispagestyle{empty}

\noindent\textcopyright\ 2013--2014 William Fang\\
{\kai 国信证券\ ·\ 厦门营业部} {\lmr\textregistered}\\[-2em]
\ \\
%\noindent All right reserved. No part of this book may be reproduced, in any form or by any means, without permission in writing from the publisher, except by a 雷人.\\
%\ \\
%\noindent The author and publisher of this book have used their best efforts in preparing this book. These efforts include the development, research, and testing of the theories, technologies and programs to determine their effectiveness. The author and publisher make no warranty of any kind, express or implied, with regard to these techniques or programs contained in this book. The author and publisher shall not be liable in any event of incidental or consequential damages in connection with, or arising out of, the furnishing, performance, or use of these techniques or programs.\\
%\ \\
%\noindent Printed in the United States of America

\begin{center}
\LARGE\enspace\bfseries{\color{themecolor} 免\ 责\ 申\ 明} \\
\base{}
\end{center}



本文仅用个人学习所用，不得以商业经营之目的进行扩散、传阅：

\begin{enumerate}
\item 未经允许不能复印／复制或向第三方详细复述本学习手册
\item All rgith reserved.
\end{enumerate}



% Open right pages: title, dedication, TOC
\makeatletter
\@openrectotrue
\makeatother



% TOC, LOF, LOT
\tableofcontents

% Open any pages: LOF, LOT, forward, preface, acknowledgements, abbreviations
\makeatletter
\@openrectofalse
\makeatother

\listoffigures
\listoftables
% \listoflcode

\renewcommand{\chaptermark}[1]{\markboth{#1}{}}
\onehalfspacing
\let\oldparskip\parskip
\addtolength{\parskip}{3pt}
\addtolength{\abovecaptionskip}{-3pt}
\addtolength{\belowcaptionskip}{-3pt}
\setlength{\parindent}{2em}

\chapter{摘\hspace{2em}要}

厦门怡惟服饰股份有限公司成为海峡股份交易中心首批挂牌企业，怡惟一直演绎现代都市女性的时尚、优雅、知性，深受中国女性消费者喜爱。这次参与挂牌，对于怡惟来说是一个新的起点。

%\chapter{再版序}

转眼间 lnotes 已经两岁了。在这两年间，它在雷友们的帮助下改正了许多缺点并茁壮地成长着。随着 \LaTeX 技术的进步，包老师感觉到零敲碎打缝缝补补已经不适应新时期的发展；lnotes 需要洗心革面重新做人，以求得人民群众的谅解。因为太懒，这个动手改版的日期被一再推迟。然而亡羊补牢，亦未为迟。

\section*{新版改进}
新版本力求实现以下几个目标：

\begin{compactenum}
    \item 系统性；结构完整，脉络清晰。
    \item 层次性；详略得当，重点突出。
    \item 进步性；技术先进，内容创新。
    \item 一致性；前后呼应，风格统一。
\end{compactenum}

古人云：知易行难。古人又云：法乎其上，得乎其中。为实现这些目标，本文作出以下具体调整：

\begin{compactenum}
    \item 全文按出版业惯例分为三大部分：前置部分，包括封面、标题、版权、献辞、目录、序、致谢等；主体部分包括九章；后置部分包括跋、附录、索引等。
    \item 第一章简介，历史回顾部分有较大扩充。
    \item 第二章入门，结构有所调整，重组了对齐和间距、特殊段落、列表等节，内容也有增强。
    \item 第三章字体，介绍电脑字体相关概念，字符集、编码、字体格式等，以及字体在 \LaTeX 中的应用。由于 \XeTeX 日趋成熟，新字体技术方案得到广泛应用；从前的一些中文解决方案和字体安装配置方法已经过时。本章由原第七章中文和第八章字体压缩合并而来，增加了 \XeTeX 相关内容。
    \item 第四章数学，结构基本不变，补充了部分示例的代码。
    \item 第五章插图，介绍图形格式及其转换，怎样插入图形，矢量绘图。增加了色彩模型介绍；若干小节标题略作调整，更加一致。
    \item \MP, PSTricks, PGF等绘图工具的介绍因篇幅较长，拆分为三个独立章节，内容亦有所增强和扩充。
    \item 第九章表格，增加了数字表格和宽表格两节，宽度控制和彩色表格有所增加和扩充。
    \item 第十章结构，介绍文档结构，标题、目录、长文档、参考文献、索引、超链接等。
    \item 第十一章布局，介绍页面尺寸，页面样式以及左右标记、章节标记等的定制，多栏、分页等。
    \item 第十二章应用，包括幻灯、书信、简历、棋谱等内容。
    \item 附录A软件和宏包。
    \item 附录B印刷简史，有图有真相。
    \item 索引，包括人物，学校、研究所、公司、政府部门等组织机构的索引。
\end{compactenum}

\section*{体例}

\begin{compactenum}
    \item 人物，西方人名正文中一般用原文，索引中加中文翻译。中国人和日本人正文中用中文，索引中加英文翻译。正文中人名第一次出现时附生卒年份。
    \item 组织、机构、公司、学校等，著名的正文中用中文，比如惠普、微软，索引中加英文。缩写，著名的正文中直接用，比如 MIT, IBM, ISO，索引中有全名；一般的先写中文名，括号内提供英文名和缩写。
    \item 人物和组织机构等，与本文没有直接关系的不加索引。
    \item 英文大小写，组织、机构、公司、学校等，一般首字母大写；其他词汇全部小写。
    \item 标点符号，中文间用全角，西文间用半角。
    \item 命令行程序、宏包、\LaTeX{}命令和环境等用等宽字体。
\end{compactenum}

%\chapter{致谢}

在本文的写作过程中，我得到了众多网友的帮助和指点，各位反动学术权威的关心和鼓励。没有你们的帮助，包老师形只影单单枪匹马马不停蹄也难以完成这件超出本人能力的事情。

在此包老师依据我国法律\footnote{《中华人民共和国感谢法》，2010年3月12日。}，首先郑重感谢党和政府的栽培，国家和人民的养育，以及有关部门的领导。感谢铁岭 TV，辽宁 TV，将来还有可能感谢 CCTV。

其次将网友们的名单公诸于众，以彰显社会良知、公民勇气。以下排名不分先后，其中多数网友来自水木清华 BBS \TeX 版和 C\TeX 论坛。

\begin{multicols}{3}
\noindent
careworn\\
Dieken\\
donated\\
gaozhch\\
garfileo\\
hkkhhk\\
Hongdong Ji\\
IMB\\
jjgod\\
Kov Chai\\
Langpku\\
LittleLeo\\
maming\\
meteorrain\\
milksea\\
northwater\\
PaladinHL\\
PiscesGold\\
primenumber\\
sesame\\
shenshou246\\
snoopyzhao\\
Sofoot\\
tex\\
Xiao Zigang\\
Xubuntu\\
yakun\\
Yiming\\
yli\\
yyzz11\\
zoho\\
张晓南\\
贾朋
\end{multicols}

也借机感谢一下4.80\footnote{166.111.4.80是清华大学校园网早期一个重要的FTP站点。}前站长 Jonny 和水木清华 BBS 前站长 Leeward。十余年前，在我开始浸淫于电脑网络时，这两位高手对我有很大的启发和影响。虽然两位高人淡出公众视野已久，但是他们为人民服务的精神却依然值得我们缅怀与尊敬。

最后还要感谢家人的理解和支持。老妻把她的博士论文给我当作学习 \LaTeX 的试验品；大女儿把她的玉照给我当作插图样板；小女儿把她的名字给我放在献辞页。

\makeatletter
\@openrectotrue  
\makeatother
\mainmatter
\renewcommand{\chaptermark}[1]{\markboth{\chaptername: #1}{}}
\renewcommand{\sectionmark}[1]{\markright{\thesection: #1}}

%%People
%\newcommand{\indexAho}{\index{people}{Aho@Alfred V. Aho, 阿尔佛雷德·艾侯}}
%\newcommand{\indexFourer}{\index{people}{Fourer@Robert Fourer, 罗伯特·福瑞尔}}
%\newcommand{\indexGay}{\index{people}{Gay@David M. Gay, 大卫·给}}
%\newcommand{\indexRitchie}{\index{people}{Ritchie@Dennis M. Ritchie, 丹尼斯·瑞奇}}
%\newcommand{\indexWeinberger}{\index{people}{Weinberger@Peter J. Weinberger, 彼得·温伯格}}
\newcommand{\indexBarclay}{\index{people}{Barclay@Robert Barclay, 罗伯特·巴克雷}}
\newcommand{\indexBernersLee}{\index{people}{Berners-Lee@Tim J. Berners-Lee, 蒂姆·伯纳斯-李}}
\newcommand{\indexBerry}{\index{people}{Berry@Karl Berry, 卡尔·白瑞}}
\newcommand{\indexCarlisle}{\index{people}{Carlisle@David P. Carlisle, 大卫·卡利斯}}
\newcommand{\indexCarlson}{\index{people}{Carlson@Chester F. Carlson, 切斯特·卡尔森}}
\newcommand{\indexCassatt}{\index{people}{Cassatt@Mary S. Cassatt, 玛丽·卡萨特}}
\newcommand{\indexCochran}{\index{people}{Cochran@Steven D. Cochran, 斯蒂文·寇克阮}}
\newcommand{\indexDaly}{\index{people}{Daly@Patrick W. Daly, 派垂克·达利}}
\newcommand{\indexDanaux}{\index{people}{Danaux@Xavier Danaux, 夏卫·丹诺克}}
\newcommand{\indexDupuis}{\index{people}{Dupuis@Étienne Dupuis, 艾提纳·杜普斯}}
\newcommand{\indexDurer}{\index{people}{Durer@Albrecht Dürer, 阿尔布雷希特·丢勒}}
\newcommand{\indexEvans}{\index{people}{Evans@David C. Evans, 大卫·埃文斯}}
\newcommand{\indexFairbairns}{\index{people}{Fairbairns@Robin Fairbairns, 罗宾·费尔贝恩}}
\newcommand{\indexFear}{\index{people}{Fear@Simon Fear, 西门怕怕}}
\newcommand{\indexFuchs}{\index{people}{Fuchs@David R. Fuchs, 大卫·福克斯}}
\newcommand{\indexGainsborough}{\index{people}{Gainsborough@Thomas Gainsborough, 托马斯·艮斯博罗}}
\newcommand{\indexGeschke}{\index{people}{Geschke@Charles M. Geschke, 查尔斯·哥谢克}}
\newcommand{\indexGirou}{\index{people}{Girou@Denis Girou, 丹尼斯·杰若}}
\newcommand{\indexGoldfarb}{\index{people}{Goldfarb@Charles F. Goldfarb, 查尔斯·歌德法布}}
\newcommand{\indexGoossens}{\index{people}{Goossens@Michel Goossen, 迈克尔·古神}}
\newcommand{\indexGoya}{\index{people}{Goya@Francisco Goya, 弗朗西斯科·戈雅}}
\newcommand{\indexGratzer}{\index{people}{Gratzer@George Grätzer, 乔治·格拉兹}}
\newcommand{\indexGutenberg}{\index{people}{Gutenberg@Johannes Gutenberg, 约翰内斯·古腾堡}}
\newcommand{\indexHagen}{\index{people}{Hagen@Hans Hagen, 汉斯·哈根}}
\newcommand{\indexHellmich}{\index{people}{Hellmich@Waldemar Hellmich, 瓦尔德玛·海米}}
\newcommand{\indexHigonnet}{\index{people}{Higonnet@René A. Higonnet, 勒内·希格涅特}}
\newcommand{\indexHobby}{\index{people}{Hobby@John D. Hobby, 约翰·爱好}}
\newcommand{\indexHoekwater}{\index{people}{Hoekwater@Taco Hoekwater, 塔克·毫克水}}
\newcommand{\indexHoover}{\index{people}{Hoover@Herbert Hoover, 赫伯特·胡佛}}
\newcommand{\indexHopfer}{\index{people}{Hopfer@Daniel Hopfer, 丹尼尔·霍卜弗}}
\newcommand{\indexJobs}{\index{people}{Jobs@Steve Jobs, 斯蒂夫·乔布斯}}
\newcommand{\indexJoy}{\index{people}{Joy@Bill N. Joy, 比尔·乔伊}}
\newcommand{\indexKastenbein}{\index{people}{Kastenbein@Karl Kastenbein, 卡尔·卡斯腾本}}
\newcommand{\indexKern}{\index{people}{Kern@Uwe Kern, 尤伟·科恩}}
\newcommand{\indexKernighan}{\index{people}{Kernighan@Brian W. Kernighan, 布莱恩·科尼汉}}
\newcommand{\indexKew}{\index{people}{Kew@Jonathan Kew, 乔纳森·酷}}
\newcommand{\indexKlic}{\index{people}{Klic@Karl Klíč, 卡尔·克里克}}
\newcommand{\indexKniagininski}{\index{people}{Kniagininski@Peter Kniagininski, 彼得·可怜人斯基}}
\newcommand{\indexKnuth}{\index{people}{Knuth@Donald E. Knuth, 高德纳}}
\newcommand{\indexLamport}{\index{people}{Lamport@Leslie Lamport, 莱斯利·兰泡特}}
\newcommand{\indexLanston}{\index{people}{Lanston@Tolbert Lanston, 托伯特·兰斯顿}}
\newcommand{\indexLemberg}{\index{people}{Lemberg@Werner Lemberg, 维尔纳·林伯格}}
\newcommand{\indexLesenko}{\index{people}{Lesenko@Sergey Lesenko, 瑟基·莱森科}}
\newcommand{\indexLevy}{\index{people}{Levy@Silvio Levy, 斯维·勒维}}
\newcommand{\indexLichtenberg}{\index{people}{Lichtenberg@Georg C. Lichtenberg, 乔治·里腾堡}}
\newcommand{\indexLorie}{\index{people}{Lorie@Raymond Lorie, 雷蒙德·劳里}}
\newcommand{\indexMarkey}{\index{people}{Markey@Nicolas Markey, 尼古拉斯·马克}}
\newcommand{\indexMcIlroy}{\index{people}{McIlroy@Malcolm D. McIlroy, 马尔考姆·麦克罗伊}}
\newcommand{\indexMerayHorvath}{\index{people}{Meray-Horvath@C. Méray Horváth, 莫雷-霍瓦斯}}
\newcommand{\indexMergenthaler}{\index{people}{Mergenthaler@Ottmar Mergenthaler, 奥特玛·摩根泰勒}}
\newcommand{\indexMiletic}{\index{people}{Miletic@Vedran Miletić, 维特恩·米尔提克}}
\newcommand{\indexMittelbach}{\index{people}{Mittelbach@Frank Mittelbach, 弗兰克·米特巴赫}}
\newcommand{\indexMorris}{\index{people}{Morris@Robert H. Morris, 罗伯特·莫里斯}}
\newcommand{\indexMoses}{\index{people}{Moses@Brooks Moses, 布鲁克斯·摩西}}
\newcommand{\indexMosher}{\index{people}{Mosher@Edward Mosher, 爱德华·莫谢}}
\newcommand{\indexMoyroud}{\index{people}{Moyroud@Louis M. Moyroud, 路易·毛绒的}}
\newcommand{\indexNewell}{\index{people}{Newell@Martin Newell, 马丁·牛维尔}}
\newcommand{\indexOberdiek}{\index{people}{Oberdiek@Heiko Oberdiek, 黑口·奥本迪克}}
\newcommand{\indexOetiker}{\index{people}{Oetiker@Tobias Oetiker, 托比·欧提克}}
\newcommand{\indexOostrum}{\index{people}{van Oostrum@Piet van Oostrum, 皮特·范·奥斯图姆}}
\newcommand{\indexOssanna}{\index{people}{Ossanna@Joseph F. Ossanna, 约瑟夫·奥萨纳}}
\newcommand{\indexOstwald}{\index{people}{Ostwald@Wilhelm Ostwald, 威廉·奥斯特瓦尔德}}
\newcommand{\indexPakin}{\index{people}{Pakin@Scott Pakin, 斯考特·培钦}}
\newcommand{\indexPatashnik}{\index{people}{Patashnik@Oren Patashnik, 奥润·帕塔尼克}}
\newcommand{\indexPorstmann}{\index{people}{Porstmann@Walter Porstmann，瓦尔特·泡茨曼}}
\newcommand{\indexRahtz}{\index{people}{Rahtz@Sebastian P. Rahtz, 塞巴斯蒂安·拉兹}}
\newcommand{\indexReagon}{\index{people}{Reagon@Ronald Reagan, 罗纳德·里根}}
\newcommand{\indexReckdahl}{\index{people}{Reckdahl@Keith Reckdahl, 凯斯·瑞克道}}
\newcommand{\indexReid}{\index{people}{Reid@Brian K. Reid, 布莱恩·瑞得}}
\newcommand{\indexRobertson}{\index{people}{Robertson@Will Robertson, 威尔·罗伯特森}}
\newcommand{\indexRochester}{\index{people}{Rochester@Wayne A. Rochester, 韦恩·罗彻斯特}}
\newcommand{\indexRokicki}{\index{people}{Rokicki@Tomas Rokicki, 托马斯·若奇奇}}
\newcommand{\indexRubel}{\index{people}{Rubel@Ira W. Rubel, 伊雷·鲁卜}}
\newcommand{\indexSaltzer}{\index{people}{Saltzer@Jerome H. Saltzer, 杰若米·萨尔兹}}
\newcommand{\indexSchandl}{\index{people}{Schandl@Bernd Schandl, 波得·辛得}}
\newcommand{\indexSchopf}{\index{people}{Schopf@Rainer Schöpf, 雨人·肖普弗}}
\newcommand{\indexSenefelder}{\index{people}{Senefelder@Alois Senefelder, 阿劳斯·塞尼菲尔德}}
\newcommand{\indexShamos}{\index{people}{Shamos@Michael I. Shamos, 迈克尔·沙毛斯}}
\newcommand{\indexSommerfeldt}{\index{people}{Sommerfeldt@Axel Sommerfeldt, 艾克赛·萨摩非特}}
\newcommand{\indexSpivak}{\index{people}{Spivak@Michael D. Spivak, 迈克尔·斯匹维克}}
\newcommand{\indexSproull}{\index{people}{Sproull@Robert F. Sproull, 罗伯特·斯普罗}}
\newcommand{\indexStallman}{\index{people}{Stallman@Richard M. Stallman, 理查德·斯托曼}}
\newcommand{\indexStarkweather}{\index{people}{Starkweather@Gary K. Starkweather, 盖瑞·斯塔克维}}
\newcommand{\indexSutherland}{\index{people}{Sutherland@Ivan E. Sutherland, 伊凡·苏泽兰}}
\newcommand{\indexTantau}{\index{people}{Tantau@Till Tantau, 提·谭陶}}
\newcommand{\indexTaylor}{\index{people}{Taylor@Philip Taylor, 菲利普·泰勒}}
\newcommand{\indexThompson}{\index{people}{Thompson@Paul A. Thompson, 保罗·汤普森}}
\newcommand{\indexUry}{\index{people}{Ury@Lesser Ury, 莱瑟·尤里}}
\newcommand{\indexVanZandt}{\index{people}{van Zandt@Timothy van Zandt, 提摩西·范·赞特}}
\newcommand{\indexVeytsman}{\index{people}{Veytsman@Boris Veytsman, 波瑞斯·韦茨曼}}
\newcommand{\indexVonSiegen}{\index{people}{von Siegen@Ludwig von Siegen, 路德维希·冯·希根}}
\newcommand{\indexVoss}{\index{people}{Voss@Herbert Voß, 赫伯特·沃斯}}
\newcommand{\indexWarnock}{\index{people}{Warnock@John E. Warnock, 约翰·沃诺克}}
\newcommand{\indexWicks}{\index{people}{Wicks@Mark A. Wicks, 马克·维克斯}}
\newcommand{\indexWilson}{\index{people}{Wilson@Peter R. Wilson，彼得·威尔森}}
\newcommand{\indexWright}{\index{people}{Wright@Joseph Wright, 约瑟夫·莱特}}
\newcommand{\indexZlatuska}{\index{people}{Zlatuska@Jiří Zlatuška, 乔治·兹拉图斯卡}}

\newcommand{\indexBiSheng}{\index{people}{毕升}}
\newcommand{\indexCaiQiwei}{\index{people}{蔡奇伟}}
\newcommand{\indexCho}{\index{people}{赵珍焕, Jin-Hwan Cho}}
\newcommand{\indexHan}{\index{people}{韩世城, {\lmr Hàn Thế Thành}}}
\newcommand{\indexHirata}{\index{people}{平田俊作, Shunsaku Hirata}}
\newcommand{\indexLee}{\index{people}{李果正}}
\newcommand{\indexLiShujun}{\index{people}{李树钧}}
\newcommand{\indexSunWenchang}{\index{people}{孙文昌}}
\newcommand{\indexUmeki}{\index{people}{梅木秀雄, Hideo Umeki}}
\newcommand{\indexWangZhen}{\index{people}{王祯}}
\newcommand{\indexWongHongling}{\index{people}{翁鸿翎}}
\newcommand{\indexWuConghui}{\index{people}{吴聪慧}}
\newcommand{\indexWuCongmin}{\index{people}{吴聪敏}}
\newcommand{\indexWuLingyun}{\index{people}{吴凌云}}
\newcommand{\indexZhangLinbo}{\index{people}{张林波}}

%Companies
%\newcommand{\indexBendix}{\index{org}{com.bendix@Bendix Corp., 本迪克斯公司}}
%\newcommand{\indexEDS}{\index{org}{com.eds@Electronic Data Systems, EDS, 电子数据系统公司}}
%\newcommand{\indexElsevier}{\index{org}{com.elsevier@Elsevier, 爱思维尔出版公司}}
%\newcommand{\indexElvenkind}{\index{org}{com.elvenkind@Elvenkind, 精灵公司}}
%\newcommand{\indexGoogle}{\index{org}{com.google@Google, 谷歌}}
%\newcommand{\indexInterleaf}{\index{org}{com.interleaf@Interleaf, 因特里弗}}
%\newcommand{\indexIntuit}{\index{org}{com.intuit@Intuit, 直觉公司}}
%\newcommand{\indexKluwer}{\index{org}{com.kluwer@Kluwer Academic Publishers, 克鲁沃学术出版公司}}
%\newcommand{\indexMCA}{\index{org}{com.mca@Massachusetts Computer Associates, 麻省计算机同伙公司}}
%\newcommand{\indexNovell}{\index{org}{com.novell@Novell, 诺维尔}}
%\newcommand{\indexRiverValley}{\index{org}{com.river@River Valley Technologies, 河谷科技}}
%\newcommand{\indexSun}{\index{org}{com.sun@Sun Microsystems, 太阳电脑}}
%\newcommand{\indexSutherlandSproull}{\index{org}{com.sutherland@Sutherland, Sproull and Associates, 苏泽兰·斯普罗及其同伙}}
\newcommand{\indexAdobe}{\index{org}{com.adobe@Adobe Systems, 土坯电脑}}
\newcommand{\indexAlphanumeric}{\index{org}{com.alphanumeric@Alphanumeric Inc., 字母数字公司}}
\newcommand{\indexApple}{\index{org}{com.apple@Apple Inc., 苹果公司}}
\newcommand{\indexATT}{\index{org}{com.att@American Telephone \& Telegraph, AT\&T, 美国电报电话公司}}
\newcommand{\indexBell}{\index{org}{com.bell@Bell Labs, 贝尔实验室}}
\newcommand{\indexCalcomp}{\index{org}{com.calcomp@Calcomp Technology, 加州电脑技术公司}}
\newcommand{\indexDEC}{\index{org}{com.dec@Digital Equipment Corp., DEC, 数字设备公司}}
\newcommand{\indexDupont}{\index{org}{com.dupont@DuPont, 杜邦}}
\newcommand{\indexEvansSutherland}{\index{org}{com.evans@Evans \& Sutherland, 埃文斯·苏泽兰公司}}
\newcommand{\indexFairchild}{\index{org}{com.fairchild@Fairchild Camera and Instrument, 仙童相机仪器}}
\newcommand{\indexFargo}{\index{org}{com.hid@Fargo Electronics, 法戈电子}}
\newcommand{\indexHal}{\index{org}{com.hal@Hal Computer Systems, 豪电脑公司}}
\newcommand{\indexHaloid}{\index{org}{com.xerox@Haloid Photographic Company, 哈罗德公司}}
\newcommand{\indexHP}{\index{org}{com.hp@Hewlett-Packard, HP, 惠普}}
\newcommand{\indexIBM}{\index{org}{com.ibm@International Business Machine, IBM, 国际商用机器公司}}
\newcommand{\indexInternationaltype}{\index{org}{com.intertype@International Typesetting Machine Company, 国际排版机公司}}
\newcommand{\indexIntertype}{\index{org}{com.intertype@Intertype Corp., 因特铸排机公司}}
\newcommand{\indexLexmark}{\index{org}{com.lexmark@Lexmark, 利盟}}
\newcommand{\indexLinotype}{\index{org}{com.linotype@Mergenthaler Linotype Company, 摩根泰勒林诺铸排机公司}}
\newcommand{\indexLucent}{\index{org}{com.lucent@Lucent Technologies, 朗讯科技}}
\newcommand{\indexMonotype}{\index{org}{com.monotype@Monotype Corp., 莫诺公司}}
\newcommand{\indexMonotypeImaging}{\index{org}{com.monotype@Monotype Imaging, 莫诺图像}}
\newcommand{\indexMonotypeLanston}{\index{org}{com.monotype@Lanston Monotype Machine Company, 兰斯顿莫诺铸排机公司}}
\newcommand{\indexMSFT}{\index{org}{com.microsoft@Microsoft, 微软}}
\newcommand{\indexOreilly}{\index{org}{com.oreilly@O'Reilly Media, 欧莱利公司}}
\newcommand{\indexPragma}{\index{org}{com.pragma@Pragma, 普瑞格玛}}
\newcommand{\indexPrintware}{\index{org}{com.printware@Printware LLC, 打印设备公司}}
\newcommand{\indexSiemens}{\index{org}{com.siemens@Siemens AG, 西门子}}
\newcommand{\indexUnilogic}{\index{org}{com.unilogic@Unilogic, 统一逻辑公司}}
\newcommand{\indexWangLabs}{\index{org}{com.wang@Wang Laboratories, 王安公司}}
\newcommand{\indexXerox}{\index{org}{com.xerox@Xerox, 施乐}}

\newcommand{\indexCanon}{\index{org}{佳能株式会社, Canon Inc.}}
\newcommand{\indexCasio}{\index{org}{卡西欧计算机株式会社, Casio Computer}}
\newcommand{\indexEpson}{\index{org}{精工爱普生株式会社, Seiko Epson}}
\newcommand{\indexOki}{\index{org}{冲电气工业株式会社, Oki Electric Industry}}
\newcommand{\indexSATO}{\index{org}{佐藤工业株式会社, SATO Corp.}}

%Universities
%\newcommand{\indexAmerican}{\index{org}{edu.american@American University, 美国大学}}
%\newcommand{\indexBrandeis}{\index{org}{edu.brandeis@Brandeis University, 布兰迪斯大学}}
%\newcommand{\indexCalTech}{\index{org}{edu.caltech@California Institute of Technology, Caltech, 加州理工学院}}
%\newcommand{\indexCase}{\index{org}{edu.case@Case Institute of Technology, 凯斯工学院}}
%\newcommand{\indexCMU}{\index{org}{edu.carnegie@Carnegie Mellon University, 卡耐基梅隆大学}}
%\newcommand{\indexColumbia}{\index{org}{edu.columbia@Columbia University, 哥伦比亚大学}}
%\newcommand{\indexCornell}{\index{org}{edu.cornell@Cornell University, 康奈尔大学}}
%\newcommand{\indexDuquesne}{\index{org}{edu.duquesne@Duquesne University, 杜肯大学}}
%\newcommand{\indexHarvard}{\index{org}{edu.harvard@Harvard University, 哈佛大学}}
%\newcommand{\indexIHEP}{\index{org}{org.ihep@Institute for High Energy Physics, 俄罗斯高能物理研究所}}
%\newcommand{\indexIMPA}{\index{org}{edu.impa, IMPA, 巴西理论数学与应用数学研究所}}
%\newcommand{\indexKettering}{\index{org}{edu.kettering@Kettering University, 凯特林大学}}
%\newcommand{\indexKIAS}{\index{org}{edu.kias@Korea Institute for Advanced Study}}
%\newcommand{\indexLondon}{\index{org}{edu.london@University of London, 伦敦大学}}
%\newcommand{\indexMainz}{\index{org}{edu.Johannes@Johannes Gutenberg University of Mainz, 约翰内斯·古腾堡美因茨大学}}
%\newcommand{\indexMaryland}{\index{org}{edu.maryland@University of Maryland, 马里兰大学}}
%\newcommand{\indexMasaryk}{\index{org}{edu.Masaryk@Masaryk University, 捷克马萨里克大学}}
%\newcommand{\indexMassachusetts}{\index{org}{edu.Massachusetts@University of Massachusetts, 马萨诸塞大学}}
%\newcommand{\indexMPS}{\index{org}{edu.mps@Max Planck Institute for Solar System Research, 马克斯·普朗克太阳系研究所}}
%\newcommand{\indexMSRI}{\index{org}{edu.msri@Mathematical Sciences Research Institute, MSRI, 数学科学研究所}}
%\newcommand{\indexOsaka}{\index{org}{edu.osaka@Osaka City University, 大阪市立大学}}
%\newcommand{\indexOxford}{\index{org}{edu.oxford@University of Oxford, 牛津大学}}
%\newcommand{\indexPrinceton}{\index{org}{edu.princeton@Princeton University, 普林斯顿大学}}
%\newcommand{\indexSouthampton}{\index{org}{edu.southampton@University of Southampton}}
%\newcommand{\indexSuwon}{\index{org}{edu.suwon@University of Suwon, 水原大学}}
%\newcommand{\indexTAMU}{\index{org}{edu.tamu@Texas A\&M University, 得克萨斯农机大学}}
%\newcommand{\indexToronto}{\index{org}{edu.toronto@University of Toronto, 多伦多大学}}
%\newcommand{\indexUMN}{\index{org}{edu.minnesota@University of Minnesota, 明尼苏达大学}}
%\newcommand{\indexUtah}{\index{org}{edu.utah@University of Utah, 犹他大学}}
%\newcommand{\indexVassar}{\index{org}{edu.vassar@Vassar College, 瓦沙学院}}
%\newcommand{\indexWayne}{\index{org}{edu.wayne@Wayne State Univesrity, 韦恩州立大学}}
%\newcommand{\indexXavier}{\index{org}{edu.xavier@Xavier University, 泽维尔大学}}
%\newcommand{\indexYale}{\index{org}{edu.yale@Yale University, 耶鲁大学}}
\newcommand{\indexMIT}{\index{org}{edu.mit@Massachusetts Institute of Technology, MIT, 麻省理工学院}}
\newcommand{\indexStanford}{\index{org}{edu.stanford@Stanford University, 斯坦福大学}}

%\newcommand{\indexHKU}{\index{org}{中国@香港大学, University of Hong Kong}}
%\newcommand{\indexCAS}{\index{org}{中国@中国科学院, Chinese Academy of Sciences, CAS}}

%Government Agencies
\newcommand{\indexDARPA}{\index{org}{gov.darpa@Defense Advanced Research Projects Agency，DARPA, 国防部高等研究计划局}}
\newcommand{\indexEC}{\index{org}{gov.eu@European Commission}}
%\newcommand{\indexNSA}{\index{org}{gov.nsa@National Security Agency, NSA, 国家安全局}}

%Organizations
%\newcommand{\indexFSF}{\index{org}{org.fsf@Free Software Foundation, FSF, 自由软件基金会}}
%\newcommand{\indexSRI}{\index{org}{org.sri@Stanford Research Institute, SRI, 斯坦福研究所}}
\newcommand{\indexAMS}{\index{org}{org.ams@American Mathematical Society, AMS, 美国数学学会}}
\newcommand{\indexANSI}{\index{org}{org.ansi@American National Standards Institute, ANSI, 美国国家标准协会}}
\newcommand{\indexCERN}{\index{org}{org.cern@European Organization for Nuclear Research, CERN, 欧洲核子研究中心}}
\newcommand{\indexECMA}{\index{org}{org.ecma@Ecma International, Ecma国际}}
\newcommand{\indexIEEE}{\index{org}{org.ieee@Institute of Electrical and Electronics Engineers, 电气电子工程师学会}}
\newcommand{\indexGNU}{\index{org}{org.gnu@GNU, 哥奴}}
\newcommand{\indexIETF}{\index{org}{org.ietf@Internet Engineering Task Force, IETF, 互联网工程任务组}}
\newcommand{\indexIMC}{\index{org}{org.imc@Internet Mail Consortium, IMC, 互联网电子邮件联盟}}
\newcommand{\indexISO}{\index{org}{org.iso@International Organization for Standardization, ISO, 国际标准化组织}}
\newcommand{\indexSIL}{\index{org}{org.sil@SIL International, SIL国际}}
\newcommand{\indexTUG}{\index{org}{org.tug@TeX User Groups, TeX用户组}}
\newcommand{\indexUnicode}{\index{org}{org.unicode@The Unicode Consortium, 统一码联盟}}
\newcommand{\indexWWWC}{\index{org}{org.w3c@World Wide Web Consortium, W3C, 万维网联盟}}

% !Mode:: "TeX:UTF-8"

\chapter{公司概况}

\section{公司基本情况}
{\bfseries 厦门怡惟服饰股份有限公司}（以下简称『厦门怡惟』）坐落于厦门市湖里区，系台湾怡惟服装有限公司的大陆独家合资子公司，首期投资超过一千万元，全权运营怡惟（EVEI）时尚女装品牌。

台湾怡惟服饰有限公司总部位于台湾高雄，自2003年成立以来，发展成为一家集研发、生产、营销、物流、信息一体化的品牌女装公司。秉承『质量先行，顾客至上，追求创新，双赢互惠』品牌经营理念，专业化的研发设计，规模化的生产物流，建立起合理的分销体系、市场网络和全方位的营销服务体系。为了向顾客提供满意的服饰产品，怡惟不惜重金力邀国内外行业内高精尖人物——新锐时尚大师、品牌运营大师等专业人士携手，共筑怡惟品牌的运营。通过8年的市场历练，精准的定位、时尚的款式、精良的工艺，EVEI时尚女装正被愈来愈多的都市女性所青睐，如今，销售门店已近200家，销售区域遍及国内各地大中城市，公司准确市场定位，发展迅速，致力打造中国强势时尚品牌女装。2011年，公司被厦门市政府评为『成长型企业』并给予政策扶持。

怡惟品牌主要面向25-40岁之间、月收入在3000以上的成熟、时尚、知性的二、三级都市女性。时尚女装是享受生活的推动力，怡惟（EVEI）致力于“Evoke：Vogue  Elegance  Intellectuality”——『引领：时尚、优雅、 知性』的设计理念，演绎着现代都市女性的时尚、优雅、知性，深受中国女性消费者喜爱。

% \smalltodol{信息来源：厦门怡惟公司网站}{厦门怡维}

厦门怡维服饰作为福建知名品牌女装，主营业务涉及品牌女装的设计、生产及销售。公司注册资本 2 245.5万元，法人代表为吴雪颖。目前公司以联营为主要合作模式，经营网点遍布福建、四川、重庆、广西、湖南、广东、江西、陕西、江苏、天津等省份及直辖市。入驻了淘宝商城、淘宝C店、京东商城、一号店、卓越亚马逊、唯品会等网络销售平台，并取得卓越的销售业绩。2011年度至2012年度，公司连续2次荣获了『厦门市成长型企业』的荣誉称号。2012年，公司品牌被中国服装协会评为『2012年度最具成长性品牌』，受到厦门市政府优惠政策的大力扶植。

\renewcommand*{\arraystretch}{1}
\setlength{\tabcolsep}{8pt}
\begin{longtable}{>{\footnotesize}r>{\footnotesize}p{10cm}}
\caption[公司概况]{厦门怡惟公司概况}\\  % 1
\hline\hline
\rowcolor{mycyan} \bfseries 项目 	& \bfseries 说明 \\  \endfirsthead          % 2 \endfirsthead          % 2
\caption[]{厦门怡惟公司概况（续表）} \\ 
\hline\hline
\rowcolor{mycyan} \bfseries 项目 	& \bfseries 说明 \\ \endhead         % 2 \endhead                % 4
\hline
\endfoot
\hline   % 内容开始
公司名称    &   厦门怡惟服饰股份有限公司 \\
英文名称	&	Xiamen Evei Garment Co., Ltd.	\\
公司类型	&	股份有限公司（自然人控股、非上市）\\
注册资本	&	 2 245.5万元 \\
法人代表	&	吴雪颖	\\
信息披露	&	滕明涛	\\
所属行业	&	纺织服装、服饰业	\\
注册地址	&	厦门市湖里区湖里大道54号第八层S1单元	\\
办公地址	&  厦门市湖里区湖里大道54号第八层S1单元	\\
经营范围	& 1、生产、加工、销售：服装、鞋帽、皮包；\newline 2、批发、零售：针纺织品、日用品、床上用品、办公用品、酒店用品、家居用品、装饰品。\\
主营业务	&	品牌女装的设计、生产及销售	\\
\bottomrule
\end{longtable}


\subsection{股本结构}
截止2013年07月18日，公司前十大股东的股本结构如下： %\smalltodol{厦门市同邦资产管理有限公司，成立于 2006 年，注册资本 5000 万元，是一家专注于挂牌上柜保荐、股权融资、项目融资、财务顾问、直接投资的综合型非银行金融服务机构。}{公司前十大股东的股本结构如下：}
%%------------------------------------------------------------------------
\renewcommand*{\arraystretch}{1}
\setlength{\tabcolsep}{8pt}
\begin{longtable}{>{\footnotesize}c>{\footnotesize}c>{\footnotesize}r>{\footnotesize}r>{\footnotesize}c}
\caption[公司前十大股东]{公司前十大股东持股情况}\\  % 1
\hline\hline
\rowcolor{mycyan} \bfseries 股东名称	&	\bfseries 股东性质	&	\bfseries 持股数量（股）	&	\bfseries 持股比例（\%）	&	\bfseries 交易是否受限	\\  \endfirsthead          % 2 \endfirsthead          % 2
\caption[]{公司前十大股东持股情况（续表）} \\ 
\hline\hline
\rowcolor{mycyan} \bfseries 股东名称	&	\bfseries 股东性质	&	\bfseries 持股数量（股）	&	\bfseries 持股比例（\%）	&	\bfseries 交易是否受限	\\   \endhead         % 2 \endhead                % 4
\hline
\endfoot
\hline   % 内容开始
吴雪颖	&	国内自然人	&	765万	&	34.06	&	是	\\
翁良庆	&	国内自然人	&	735万	&	32.73	&	是	\\
蔡亚潮	&	国内自然人	&	150万	&	6.68	&	否	\\
翁立芳	&	国内自然人	&	100万	&	4.45	&	否	\\
厦门同邦	&	其他组织	&	75万	&	3.34	&	否	\\
杨美香	&	国内自然人	&	50万	&	2.22	&	否	\\
曾桑潭	&	国内自然人	&	50万	&	2.22	&	否	\\
吴春红	&	国内自然人	&	40万	&	1.78	&	否	\\
滕明涛	&	国内自然人	&	30万	&	1.34	&	否	\\
谢翠玲	&	国内自然人	&	30万	&	1.34	&	否	\\
\midrule
合计	&	——	&	2 025万	&	90.18	&		—— 	\\
\bottomrule
\end{longtable}
\begin{note}
其中，厦门市同邦资产管理有限公司作为机构管理入股，成立于2006年，注册资本5000万元，是一家专注于挂牌上柜保荐、股权融资、项目融资、财务顾问、直接投资的综合型非银行金融服务机构。
\end{note}
%%------------------------------------------------------------------------
\subsection{组织结构}
目前，公司法定代表人为吴雪颖。公司的管理层组织架构为
%%------------------------------------------------------------------------
\renewcommand*{\arraystretch}{1}
\setlength{\tabcolsep}{8pt}
\begin{longtable}{>{\footnotesize}c>{\footnotesize}l>{\footnotesize}c>{\footnotesize}c>{\footnotesize}r>{\footnotesize}r>{\footnotesize}r}
\caption[公司主要管理者]{公司主要管理者介绍}\\  % 1
\hline\hline
\rowcolor{mycyan} \bfseries 高管名称	&	\bfseries 职位	&	\bfseries 出生年份	& \bfseries	性别	& \bfseries	学历	&	\bfseries 所持股份（股）	& \bfseries	持股比例（\%）	\\  \endfirsthead          % 2 \endfirsthead          % 2
\caption[]{公司主要管理者介绍（续表）} \\ 
\hline\hline
\rowcolor{mycyan} \bfseries 高管名称	&	\bfseries 职位	&	\bfseries 出生年份	& \bfseries	性别	& \bfseries	学历	&	\bfseries 所持股份（股）	& \bfseries	持股比例（\%）	\\   \endhead         % 2 \endhead                % 4
\hline
\endfoot
\hline   % 内容开始
吴雪颖	&	董事长	&	1973年	&	女	&	本科	&	765万	&	34.06	\\
蔡亚潮	&	副董事长	&	1968年	&	男	&	本科	&	150万	&	6.68	\\
翁立芳	&	总经理兼董事	&	1975年	&	男	&	大专	&	100万	&	4.45	\\
匡全军	&	总经理兼董事	&	1979年	&	男	&	大专	&	无	&	——	\\
郑\hspace{2ex}桦	&	董事	&	1969年	&	女	&	本科	&	30万	&	1.34	\\
曾根福	&	董事	&	1963年	&	男	&	本科	&	10万	&	0.46	\\
滕明涛	&	董事会秘书	&	1969年	&	男	&	本科	&	30万	&	1.34	\\
林明光	&	董事	&	1978年	&	男	&	中专	&	无	&	——	\\
翁良庆	&	监事会主席	&	1950年	&	男	&	初中	&	735万	&	32.73	\\
吴永会	&	监事	&	1979年	&	男	&	大专	&	无	&	——	\\
兰\hspace{2ex}丽	&	监事	&	1986年	&	女	&	大专	&	无	&	——	\\
陈群英	&	财务总监	&	1971年	&	女	&	本科	&	无	&	——	\\
\bottomrule
\end{longtable}
%%------------------------------------------------------------------------

%%------------------------------------------------------------------------
\subsection{主要管理者介绍}
厦门怡惟公司现有管理层，其具体情况如下。\\
\vspace{1ex}

\begin{mdfbox}[吴雪颖]
\end{mdfbox}

从怡惟公司主要管理层成员结构上看
\section{公司经营与财务情况}
%%------------------------------------------------------------------------
\subsection{经营情况介绍}
厦门怡惟公司的主营业务涉及品牌女装的设计、生产及销售，来自主营业务的收入在过去的两个会计年度间表现出强劲的增长态势。公司注重品牌的建设及推广，采用以联营为主，自营、电商、加盟相结合的商业运营模式，目前已将业务扩展至福建、四川等多个省份及直辖市，迅速取得了一定的市场份额。同时，公司也在积极开拓新型销售渠道，建立、完善了一套多层次、一体化的完整销售体系，不仅充分发挥实体店面客户亲和力、品牌忠诚感、口碑知名度的传统渠道的营销能力，还积极拓展网络销售渠道，入驻了淘宝商城、淘宝C店、京东商城、一号店、卓越亚马逊、唯品会等网络销售平台，并取得卓越的销售业绩。

近几年来怡惟公司在资产规模、市场占有率、盈利能力、净利润等方面增长显著，正逐步跻身为女装服饰行业的第一梯队。公司整体资产总额保持高速增长，在2011年至2013年间，分别增长了121.62\%、50.62\%，公司的总资产报酬率在三年间分别达到了15.20\%、9.03\%与18.35\%，这主要是由于公司通过减少公司负债总额，同时通过大量的权益性资产融资，增强自身的盈利创造能力。公司具有较强的研发设计团队，目前能够较好的预测和把握服饰流行动向，其主营业务收入分别上升了40.74\%、278.08\%，也显示了公司在未来面临着十分广阔的发展空间。公司主要产品的原材料主要为各种面、辅料。公司拥有系统的采购管理办法和外包业务操作流程，与一些主要供应商和外包厂商建立了稳定良好的合作关系。此外，公司通过优化工艺流程降低了单位产品原料用量，降低了原材料价格波动给公司带来的经营压力，维持着稳定的毛利率，公司在2011年至2013年间，销售毛利率分别达到41.14	\%、37.18\%和35.66\%。


与此同时，需要提请投资者关注公司的存货资产规模增长速度。由于服饰市场具有较强的时尚性与流行性，如果公司对流行时尚和市场需求判断失误，不能及时引导时尚潮流、推出迎合时尚趋势的产品，从而不能全面满足消费者需求，则其大规模的存货积压有可能无法通过市场销售转化为公司的利润收入。产品的大量滞销与积压，将对公司未来的经营业绩产生不利影响。
%%------------------------------------------------------------------------


%%------------------------------------------------------------------------

\subsection{财务情况介绍}
根据现有企业年度会计报表资料，我们将部分重要财务指标罗列如下，具体财务分析将在下一章节分析。

从盈利能力指标分析，怡惟公司在过去的三年时间内呈现出『U型』的特征，即在2011年期间，公司来自主营业务取得的收入使得其保持较为强劲的盈利能力，权益资本投资回报率（ROE）为27.28\%，往后的一年内虽然主营业务依然发展迅速，但由于公司规模扩大导致其管理费用大幅增加，由此造成权益资本回报率减低为8.83\%，其后的2013年度则通过减低相应的管理费用，增加权益资本比重，使得该年度的权益资本回报率上升为15.44\%。同样，税前投入资本利润率分别为25.55\%、7.47\%、16.65\%。

针对公司的资金流动性与债权偿还能力的分析显示，怡惟公司具备较为稳健的偿还债务能力，公司总资产负债率持续减低，流动比率与速动比率在显著增加，这保证了公司在偿还债务方面的能力，不存在无法实现兑付偿还的风险。而公司的资金周转速度逐年加快，资产使用效率也在逐步提高，为公司未来发展提供充分的资金流动性。公司的权益资产比值从2011年的40\%上升为89\%，权益乘数达到1.13，这显示了公司股东对公司管理决策权的绝对控制权，增强了公司在重大战略决策的宏观把握。
%%------------------------------------------------------------------------
%%------------------------------------------------------------------------
\renewcommand*{\arraystretch}{0.8}
\setlength{\tabcolsep}{5pt}
\begin{longtable}{>{\footnotesize}l>{\footnotesize}r>{\footnotesize}r>{\footnotesize}r>{\footnotesize}r>{\footnotesize}r}
\caption[主要财务指标]{主要财务指标}\\  % 1
\hline\hline
\rowcolor{mycyan} {\bfseries \footnotesize  项目} & {\bfseries \footnotesize  2011}\hspace{2ex} & {\bfseries \footnotesize   2012}\hspace{2ex} &   {\bfseries \footnotesize   2013}\hspace{2ex}  &{\bfseries \footnotesize  11-12(\%)} &{\bfseries \footnotesize  12-13(\%)} \\  \endfirsthead          % 2
\caption[]{主要财务指标（续表）} \\ 
\hline\hline
\rowcolor{mycyan} {\bfseries \footnotesize  项目} & {\bfseries \footnotesize  2011}\hspace{2ex} & {\bfseries \footnotesize   2012}\hspace{2ex} &   {\bfseries \footnotesize   2013}\hspace{2ex}  &{\bfseries \footnotesize  11-12(\%)} &{\bfseries \footnotesize  12-13(\%)} \\ \endhead                % 4
\hline
\endfoot
\hline   % 内容开始
总资产	&	17,931,067.33	&	39,739,359.18	&	59,853,895.87	&	121.62	&	50.62	\\
负债总额	&	10,696,565.35	&	5,775,686.01	&	5,445,518.93	&	-46.00	&	-5.72	\\
所有者权益	&	7,234,501.98	&	33,963,673.17	&	54,349,779.31	&	369.47	&	60.02	\\
\midrule
主营收入	&	8,406,216.30	&	11,831,083.55	&	44,731,367.24	&	40.74	&	278.08	\\
营业利润	&	2,668,082.02	&	2,401,723.12	&	9,140,389.39	&	-9.98	&	280.58	\\
利润总额	&	2,595,487.43	&	2,437,480.83	&	9,089,971.70	&	-6.09	&	272.92	\\
投入资本	&	10,670,756.42	&	34,869,112.38	&	54,877,398.54	&	226.77	&	57.38	\\
EBIT	&	2,726,228.00	&	2,604,015.66	&	9,138,701.35	&	-4.48	&	250.95	\\
\midrule
销售毛利率	&	41.14	&	37.18	&	35.66	&	-9.61	&	-4.09	\\
销售利润率	&	23.47	&	15.38	&	15.24	&	-34.50	&	-0.88	\\
成本费用利润率	&	47.55	&	25.87	&	25.70	&	-45.59	&	-0.66	\\
盈余现金保障倍数	&	——	&	-532.33	&	-17.60	&	——	&	-96.69	\\
总资产报酬率	&	15.20	&	9.03	&	18.35	&	-40.60	&	103.22	\\
总资产收益率	&	11.00	&	6.31	&	13.69	&	-42.67	&	117.01	\\
营业利润率	&	32.43	&	22.01	&	20.43	&	-32.13	&	-7.18	\\
\midrule
税前总资产盈利力	&	15.20	&	9.03	&	18.35	&	-40.60	&	103.22	\\
税后总资产盈利力	&	12.62	&	7.50	&	15.23	&	-40.60	&	103.22	\\
权益资本利润率(ROE)	&	27.28	&	8.83	&	15.44	&	-67.62	&	74.82	\\
税前(ROIC)	&	25.55	&	7.47	&	16.65	&	-70.77	&	122.99	\\
税后(ROIC)	&	21.21	&	6.20	&	13.82	&	-70.77	&	122.99	\\
\midrule
流动比率	&	1.50	&	5.50	&	8.01	&	267.71	&	45.62	\\
速动比率	&	1.21	&	2.44	&	3.02	&	101.62	&	23.60	\\
现金流动负债比率	&	——	&	-167.67	&	-22.03	&	——	&	-86.86	\\
已获利息倍数	&	20.85	&	15.64	&	187.54	&	-25.01	&	1,099.37	\\
营运资本需求量比率	&	29.53	&	65.37	&	63.74	&	121.34	&	-2.49	\\
总资产负债率	&	0.60	&	0.29	&	0.11	&	-52.12	&	-60.55	\\
\midrule
权益资产比	&	0.40	&	0.71	&	0.89	&	77.06	&	24.13	\\
权益乘数	&	2.48	&	1.40	&	1.13	&	-43.52	&	-19.44	\\
基于EBIT的利息保障倍数	&	20.85	&	15.64	&	187.54	&	-25.01	&	1,099.37	\\
总资产周转次数	&	0.47	&	0.41	&	0.90	&	-12.48	&	118.93	\\
固定资产周转次数	&	4.40	&	2.99	&	6.74	&	-32.07	&	125.11	\\
\bottomrule
\end{longtable}

\end{document}
